
\sourcenote{Genre: alias rule (tier 3) --- a respelling template; the
batch's shared template is Recipe 1 of the \rulename{StreettDegenerate}
section.}

\begin{brecipe}[the filled-in application --- one pair, normalized]
\label{rec:rnnf-filled}
Fix a Streett pair $(\rho, \varphi)$ inside a \rulename{QStreettBundle}
node, with $\rho$ complement-typed. Write $\mathrm{nnf}$ for the
negation-normal-form pass on the propositional skeleton (atoms,
$\neg$, $\wedge$, $\vee$): negations pushed inward by the
\rulename{NNFDuality} rewrites until only literals remain. The rewrite
is
\[
\bigl(\mathrm{cmpl}\,\rho,\ \varphi\bigr)
\quad\rightsquigarrow\quad
\bigl(\mathrm{cmpl}\,\bigl(\mathrm{elimNeg}(\mathrm{nnf}(\rho))\bigr),\ \varphi\bigr)
\]
--- the body normalized, the complement tag \emph{retained}: per the
source discipline the complement applies at the consumption boundary
(the complement-request rule), never as a literal. After the rewrite
every literal of the request body is positive (atom-swapped where it
was negative), and the pair is in exactly the input shape the
bundle and the complement-request rule expect. One pair to one pair,
inside one node: the alias shape.
\end{brecipe}

\begin{brecipe}[payload transport]\label{rec:rnnf-transport}
\begin{center}\small
\begin{tabular}{@{}ll@{}}
\textbf{un-normalized pair} & \textbf{normalized pair}\\
\midrule
request body $\rho$ ($\mathrm{cmpl}$ tag retained) & $\mathrm{elimNeg}(\mathrm{nnf}(\rho))$, negative literals eliminated through dual atoms\\
response formula, certificates, rows & verbatim\\
polarity bookkeeping & the polarity theorem (Proposition~\ref{prop:rnnf-polarity}); nothing stored\\
\end{tabular}
\end{center}
Unpack-function slot: no verified payload function in this edition (batch template); the
structural mechanization exists (the normalizer and its theorems, named
below; not hash-bound in this edition).
\end{brecipe}

\begin{bprop}[the exactness identity --- denotation preservation]
\label{prop:rnnf-exact}
$[\![\mathrm{nnf}(\rho)]\!] = [\![\rho]\!]$, and eliminating negative
literals through dual atoms preserves the denotation.
\end{bprop}
\begin{proof}
Induction on $\rho$, with the \rulename{NNFDuality} equalities as the
cases --- double negation and De Morgan drive the recursion, the atom
swaps close the leaves. The elimination step substitutes, for each
negative literal, the dual atom denoting the complement; equality is
preserved literal by literal. (Machine-checked: \texttt{denot\_nnf} and
\texttt{denot\_elimNeg}; the eliminated form has no negative literals at
all, \texttt{negLits\_elimNeg}. Binding pending.) Complements of equal
sets are equal, so the retained $\mathrm{cmpl}$ tag's
consumption-boundary denotation $S \smallsetminus [\![\rho]\!]$ is
preserved as well --- the rewrite itself never discharges the
complement.
\end{proof}

\begin{bprop}[the polarity theorem]\label{prop:rnnf-polarity}
The surviving negative literals of $\mathrm{nnf}(\rho)$ are exactly the
atom occurrences of $\rho$ at odd polarity (an odd number of enclosing
negations). Consequently the normal form is free of negated occurrences
of a given atom class if and only if the input has no odd-polarity
occurrence of that class.
\end{bprop}
\begin{proof}
Strengthen to the polarity-indexed pass (normalizing $\rho$ or
$\neg\rho$ according to a polarity bit) and induct: at an atom the bit
decides the literal's sign; at a negation the bit flips; at $\wedge$ and
$\vee$ the bit distributes (De Morgan swapping the connective when
negative). The negative-literal set on one side and the odd-polarity
occurrence set on the other satisfy the same recursion. The rejection
face is a rewrite of the same equality. (Machine-checked:
\texttt{negLits\_nnf}, via the auxiliary \texttt{negLits\_nnfAux}, and
\texttt{nnf\_no\_negLit\_iff}; negation-freedom of the output holds by
\emph{type} --- the normal-form grammar has no negation constructor.
Binding pending.)
\end{proof}

\begin{bprop}[soundness corollary]\label{prop:rnnf-sound}
A bundle node whose pair is normalized concludes identically to the
un-normalized node: through \rulename{QStreettBundle} directly for
base-typed requests, and through the \rulename{StreettCmplRequest}
weakening for the consumed complement --- with
Proposition~\ref{prop:rnnf-exact} guaranteeing both nodes constrain the
same regions.
\end{bprop}
\begin{proof}
Immediate from Proposition~\ref{prop:rnnf-exact} and the target rules'
theorems; the alias adds no semantic content of its own.
\end{proof}

\begin{bremark}[what the polarity theorem is protecting]
The typing discipline rejects negated probability literals at even
surface polarity --- a negated lower bound is an upper bound, and the
certificate calculus deliberately has no rule concluding probability
upper bounds through this pass. The theorem makes the rejection
\emph{stable under normalization}: no sequence of dualities can
manufacture or hide a rejected literal, so the check may run on the
surface formula once, before normalization, with the same verdict. A
guard that had to chase literals through rewrites would be a moving
target; a polarity count is not.
\end{bremark}

\begin{anchors}
normalizer, negation-freedom by type & \texttt{Form.nnf}, \texttt{NForm} (\texttt{QtslAux.NNFDuality}; machine-checked, not hash-bound in this edition)\\
Proposition \ref{prop:rnnf-exact} & \texttt{denot\_nnf}, \texttt{denot\_elimNeg}, \texttt{negLits\_elimNeg} (same module, not hash-bound in this edition)\\
Proposition \ref{prop:rnnf-polarity} & \texttt{negLits\_nnf}, \texttt{nnf\_no\_negLit\_iff} (same module, not hash-bound in this edition)\\
Proposition \ref{prop:rnnf-sound} targets & the \rulename{QStreettBundle} and \rulename{StreettCmplRequest} sections\\
\end{anchors}
